\documentclass[10pt,
% a4paper,
%twocolumn,
fleqn,
%landscape, 
% papersize,
dvipdfmx,
uplatex
]{jsarticle}



\def\maru#1{\textcircled{\scriptsize#1}}%丸囲み番号

% \RequirePackage[2020/09/30]{platexrelease}

%太字設定
\usepackage[deluxe]{otf}

\usepackage{emathEy}

\usepackage[g]{esvect}

%大きな文字
\usepackage{fix-cm}

%定理環境
\usepackage{emathThm}
%\theoremstyle{boxed}
\theorembodyfont{\normalfont}
\newtheorem{Question}{問題}[subsection]
\newtheorem{Q}{}[subsection]
\newtheorem{question}[Question]{}
\newtheorem{quuestion}{}[subsection]

%セクション，大問番号のデザイン
\renewcommand{\labelenumi}{(\arabic{enumi})}
\renewcommand{\theenumii}{\alph{enumii})}
\renewcommand{\thesection}{第\arabic{section}章}

%用紙サイズの詳細設定
\usepackage{bxpapersize}
\papersizesetup{size={80mm,45mm}}
\usepackage[top=0.7zw,bottom=0truemm,left=3truemm,right=133truemm]{geometry}
\usepackage[dvipdfmx]{graphicx}

%余白など
\usepackage{setspace} % 行間
\setlength{\mathindent}{1zw}
\setlength\parindent{0pt}


%色カラーに関する設定
\usepackage{color}
\definecolor{shiro}{rgb}{0.95703125,0.87109375,0.7421875}
\definecolor{kin}{rgb}{0.95703125,0.87109375,0.7421875}
\definecolor{orange}{rgb}{1,0.7,0.2}
\definecolor{bradorange}{rgb}{1,0.5,0}
\definecolor{pink}{rgb}{0.9176,0.5686,0.5960}
\definecolor{mizu}{rgb}{0.6156,0.8,0.9955}
\color{kin}
% \pagecolor{hukamido}

\usepackage{at}%図の配置
% \usepackage{wallpaper}

\begin{document}



\at(0cm,0cm){\includegraphics[width=8cm,bb=0 0 1920 1080]{./youtube/thumbnails/templates/smart_background/複素数と方程式.jpeg}}
{\color{orange}\bf\boldmath\LARGE\underline{対称式の連立方程式}}\vspace{0.3zw}

\large 
\bf\boldmath 問．$x,\;y$の連立方程式

\vspace{0.3zw}
\hspace{0.5zw}$2x+2y+xy=3a-1,\;\vspace{0.3zw}$



\vspace{0.3zw}
\hspace{0.5zw}$x+y+xy=a\vspace{0.3zw}$


が$x,\;y$が実数である解をもつような実数$a$の範囲を求めよ．
\at(6.4cm,0.2cm){\small\color{bradorange}$\overset{\text{複素数と方程式}}{\text{典型}}$}


\newpage



\at(0cm,0cm){\includegraphics[width=8cm,bb=0 0 1920 1080]{./youtube/thumbnails/templates/smart_background/複素数と方程式.jpeg}}
{\color{orange}\bf\boldmath\Large\underline{解の条件$〜$解と係数の関係$〜$}}\vspace{0.3zw}

\LARGE 
\bf\boldmath 問．方程式$2x^2+ax+6=0$の$2$解のうち，一方が他方の$3$倍であるように実数$a$の値を定めよ．
\at(6.4cm,0.2cm){\small\color{bradorange}$\overset{\text{複素数と方程式}}{\text{典型}}$}


\newpage



\at(0cm,0cm){\includegraphics[width=8cm,bb=0 0 1920 1080]{./youtube/thumbnails/templates/smart_background/複素数と方程式.jpeg}}
{\color{orange}\bf\boldmath\LARGE\underline{$2$解から$2$次方程式を作る}}\vspace{0.3zw}

\LARGE 
\bf\boldmath 問．$3x^2-5x+1=0$の$2$つの解を$\alpha ,\;\beta$とするとき，$\alpha ^3,\;\beta ^3$を$2$つの解にもつ$2$次方程式を$1$つ作れ．
\at(6.4cm,0.2cm){\small\color{bradorange}$\overset{\text{複素数と方程式}}{\text{典型}}$}


\newpage



\at(0cm,0cm){\includegraphics[width=8cm,bb=0 0 1920 1080]{./youtube/thumbnails/templates/smart_background/複素数と方程式.jpeg}}
{\color{orange}\bf\boldmath\Large\underline{$2$次方程式の$2$解$\alpha ,\;\beta$の対称式}}\vspace{0.3zw}

\Large 
\bf\boldmath 問．$2x^2-3x+4=0$の$2$つの解を$\alpha ,\;\beta$とおくとき，
$\left(2-\alpha \right)\left(2-\beta \right)$の値を求めよ．
\at(6.4cm,0.2cm){\small\color{bradorange}$\overset{\text{複素数と方程式}}{\text{典型}}$}


\newpage



\at(0cm,0cm){\includegraphics[width=8cm,bb=0 0 1920 1080]{./youtube/thumbnails/templates/smart_background/複素数と方程式.jpeg}}
{\color{orange}\bf\boldmath\LARGE\underline{解と係数の関係の証明}}\vspace{0.3zw}

\large 
\bf\boldmath 問．$2$次方程式$ax^2+bx+c=0$において，次が成り立つことを示せ．
$2$解が$x=\alpha ,\;\beta$である$\iff\alpha +\beta =-\bunsuu{b}{a},\;\alpha \beta =\bunsuu{c}{a}$
\at(6.4cm,0.2cm){\small\color{bradorange}$\overset{\text{複素数と方程式}}{\text{典型}}$}


\newpage



\at(0cm,0cm){\includegraphics[width=8cm,bb=0 0 1920 1080]{./youtube/thumbnails/templates/smart_background/複素数と方程式.jpeg}}
{\color{orange}\bf\boldmath\large\underline{異なる$3$つの実数解をもつ$3$次方程式}}\vspace{0.3zw}

\Large 
\bf\boldmath 問．$3$次方程式

\vspace{0.3zw}
\hspace{0.5zw}$x^3+\left(a-1\right)x^2-\left(a-4\right)x-4=0\vspace{0.3zw}$


が異なる$3$つの実数解をもつような定数$a$の値の範囲を求めよ．
\at(6.4cm,0.2cm){\small\color{bradorange}$\overset{\text{複素数と方程式}}{\text{典型}}$}


\newpage



\at(0cm,0cm){\includegraphics[width=8cm,bb=0 0 1920 1080]{./youtube/thumbnails/templates/smart_background/複素数と方程式.jpeg}}
{\color{orange}\bf\boldmath\large\underline{解の配置$〜$解と係数の関係の利用$〜$}}\vspace{0.3zw}

\normalsize 
\bf\boldmath 問．$2$次方程式$x^2-2\left(m-2\right)x-m+{14}=0$が次のようなことなる$2$つの解を持つとき，定数$m$の値の範囲を求めよ．\\
(1)  ともに正の解\\
(2)  ともに負の解\\
(3)  符号が異なる解\\
(4)  ともに$1$より大きい\\

\at(6.4cm,0.2cm){\small\color{bradorange}$\overset{\text{複素数と方程式}}{\text{典型}}$}


\newpage



\at(0cm,0cm){\includegraphics[width=8cm,bb=0 0 1920 1080]{./youtube/thumbnails/templates/smart_background/複素数と方程式.jpeg}}
{\color{orange}\bf\boldmath\LARGE\underline{$2$重解をもつ$3$次方程式}}\vspace{0.3zw}

\LARGE 
\bf\boldmath 問．$3$次方程式$x^3+\left(a-4\right)x-2a=0$が$2$重解をもつとき，実数$a$の値を求めよ．
\at(6.4cm,0.2cm){\small\color{bradorange}$\overset{\text{複素数と方程式}}{\text{典型}}$}


\end{document}

